% Section 11.2, from lectures/L11/appendix/b_traffic_light.md.

\section{Genomarbetat exempel: från flödesplan till program}
\label{c11:app:b}

\subsection{Uppgiften: vägarbetsljuset}
\label{c11:b1}
Vid ett vägarbete är vägen smal nog för ett körfält, och trafiken från de två hållen turas om. Vid
varje ände står ett ljus med en röd och en grön lampa. Det här exemplet programmerar det ena ljuset:
\textbf{rött i 4 sekunder, grönt i 4 sekunder, och om igen, för alltid.}

Exemplet är medvetet enklare än trafikljuset i slutuppgiften: två lampor och två tillstånd i
stället för tre lampor och fyra. Metoden är densamma, och det är metoden avsnittet handlar om:
\begin{enumerate}
  \item Skriv upp \textbf{tillstånden}, vilka lampor som lyser i vart och ett, och hur länge.
  \item Rita \textbf{flödesplanen}.
  \item Välj \textbf{stift och bitmönster}.
  \item Skriv \textbf{programmet}, ett tillstånd i taget, med subrutiner för det som upprepas.
  \item \textbf{Testa} i simulatorn, och sedan på kortet.
\end{enumerate}

Lamporna är lysdioder på kortet: röd på stift 8 (PB0) och grön på stift 10 (PB2), kopplade som i
slutuppgiften.

\subsection{Tillstånden}
\label{c11:b2}
Ett styrobjekt som går igenom samma steg om och om igen beskrivs bäst som en följd av
\textbf{tillstånd}. Ett tillstånd säger vilka utgångar som är på, och hur länge systemet stannar
där:

\begin{narrowtable}{@{}lcccl@{}}
  \thead{Tillstånd} & \thead{Röd (PB0)} & \thead{Grön (PB2)} & \thead{Tid} & \thead{Därefter} \\
  \midrule
  Stopp & tänd & släckt & 4~s & Kör \\
  Kör & släckt & tänd & 4~s & Stopp \\
\end{narrowtable}

\noindent Det är samma idé som tillståndsföljden för en räknare i kapitel~\ref{c6:ch}: systemet är
i ett tillstånd i taget, och går till nästa vid en bestämd händelse. För en räknare var händelsen
en klockflank. Här är det att en viss tid har gått.

Tabellen är värd att skriva även när uppgiften är så här liten. Den är specifikationen, och det
färdiga programmet ska kunna kontrolleras mot den rad för rad.

\subsection{Flödesplanen}
\label{c11:b3}

\bookfigure{l11_roadwork_flowchart}{Flödesplanen för vägarbetsljuset.}{c11:fig:roadwork_flowchart}

\noindent Varje tillstånd blir två symboler: en utmatning som tänder rätt lampor, och ett
subrutinanrop som väntar. Loopen tillbaka har inget villkor, eftersom ljuset aldrig slutar. Det
enda som händer före loopen är initieringen: stackpekaren, eftersom programmet anropar subrutiner,
och de två stiften som utgångar.

\subsection{Subrutinen \code{delay_s}}
\label{c11:b4}
Programmet behöver vänta hela sekunder, men \code{delay_ms} från \secref{c10:c4} räcker bara till
255~ms. En ny subrutin, \code{delay_s}, tar antalet sekunder i \code{r24} och anropar
\code{delay_ms} med 250~ms fyra gånger per sekund:

\needspace{21\baselineskip}
\begin{asmcode}
;--------------------------------------------------------------------------------
; delay_s: Waits r24 seconds (1-63), as 4 x r24 calls of delay_ms with 250 ms each.
; A call takes 16 000 084 x r24 + 18 cycles: about 5 microseconds too long per second.
;--------------------------------------------------------------------------------
delay_s:
    push r24                        ; Save the registers this subroutine changes.
    push r25
    mov r25, r24
    lsl r25                         ; r25 = 2 x r24 ...
    lsl r25                         ; ... = 4 x r24: the number of quarter seconds.
    ldi r24, QUARTER_MS             ; A quarter of a second is 250 ms.
delay_s_loop:
    rcall delay_ms
    dec r25
    brne delay_s_loop
    pop r25                         ; Restore them, in the opposite order.
    pop r24
    ret
\end{asmcode}

\noindent Tre saker i den är värda att lägga märke till:
\begin{itemize}
  \item \textbf{\code{lsl} två gånger multiplicerar med fyra}, som i kapitel~\ref{c8:ch}: varje
    skift åt vänster dubblar talet. Tre sekunder blir tolv kvartar.
  \item \textbf{Den följer kursens regel.} Den ändrar \code{r24} och \code{r25}, och sparar båda.
    Efter \code{ldi r24, 4} och \code{rcall delay_s} är \code{r24} fortfarande 4, och kan användas
    igen.
  \item \textbf{Den fungerar upp till 63 sekunder.} Fyra gånger 64 är 256, som inte får plats i
    \code{r25}.
\end{itemize}

Konstanten \code{QUARTER_MS} står överst i programmet och är normalt 250:

\needspace{5\baselineskip}
\begin{asmcode}
.equ QUARTER_MS = 250               ; 250 ms. Set it to 1 to run about 250 times
                                    ; faster in the simulator; set it back before
                                    ; flashing.
\end{asmcode}

\noindent Den finns för simulatorns skull, se \secref{c11:b6}.

\textbf{Hur exakt är den?} Varje kvart är ett anrop av \code{delay_ms} med 250, alltså
\mbox{\code{16 000 · 250 + 18}} cykler, plus 3 cykler för \code{dec} och \code{brne} i loopen. Fyra
kvartar blir 16\,000\,084 cykler, och hela anropet \mbox{\code{16 000 084 · r24 + 18}} cykler. En
sekund blir alltså ungefär 5~µs för lång. Det är långt mindre än felet i kortets klocka
(\secref{c11:a7}).

\subsection{Programmet, rad för rad}
\label{c11:b5}
Hela programmet finns i \code{roadwork_light.asm}. Konstanter och initiering först:

\needspace{14\baselineskip}
\begin{asmcode}
.equ RED = 0                        ; PB0, Arduino pin 8.
.equ GREEN = 2                      ; PB2, Arduino pin 10.

main:
    ldi r16, high(RAMEND)           ; SP = RAMEND = 0x08FF.
    out SPH, r16
    ldi r16, low(RAMEND)
    out SPL, r16

    ldi r16, (1 << RED) | (1 << GREEN)
    out DDRB, r16                   ; The two lights are outputs.
\end{asmcode}

\noindent\code{(1 << RED)} är ett uttryck som assemblern räknar ut: en etta skiftad \code{RED} steg
åt vänster, alltså \code{0b00000001}, och \code{(1 << GREEN)} blir \code{0b00000100}. \code{|} är
OR, bit för bit, så båda tillsammans blir \code{0b00000101}. Skrivsättet gör att raden säger
\emph{vilka} lampor det gäller, i stället för bara ett tal, och om grön flyttas till ett annat
stift behöver bara \code{.equ}-raden ändras.

Sedan de två tillstånden, i en loop:

\needspace{14\baselineskip}
\begin{asmcode}
loop:
    ldi r16, (1 << RED)             ; Red, 4 s.
    out PORTB, r16
    ldi r24, 4
    rcall delay_s

    ldi r16, (1 << GREEN)           ; Green, 4 s.
    out PORTB, r16
    ldi r24, 4
    rcall delay_s

    rjmp loop                       ; And round again.
\end{asmcode}

\noindent Varje tillstånd är fyra rader: bitmönstret, ut på porten, tiden, vänta. Programmet läses
som tabellen i \secref{c11:b2}, rad för rad, och det är ingen slump: tabellen skrevs först.

Lägg märke till att \code{out PORTB, r16} skriver \textbf{hela} porten. Den tänder en lampa och
släcker den andra i samma instruktion, och det är därför bitmönstret för \enquote{Kör} är
\code{0b00000100} och inte bara \enquote{tänd grön}. Med \code{sbi} och \code{cbi} hade samma sak
krävt två instruktioner per tillstånd.

\subsection{Testa i simulatorn}
\label{c11:b6}
Ett tillstånd på 4 sekunder är 64 miljoner cykler, och simulatorn behöver en lång stund för att
köra dem. Därför finns \code{QUARTER_MS}: sätt den till \textbf{1} medan programmet testas i
simulatorn. Då tar varje \enquote{sekund} ungefär 4~ms, och hela cykeln syns på ett ögonblick, i
samma ordning och med samma förhållande mellan tiderna.

Så testar du:
\begin{enumerate}
  \item Sätt \code{QUARTER_MS} till 1 och bygg.
  \item Sätt en brytpunkt på varje \code{out PORTB, r16}.
  \item Kör med \textbf{Continue} (\textbf{F5}). Vid varje stopp: stega ett steg med \textbf{F11}
    och kontrollera i I/O-fönstret att rätt bitar i \code{PORTB} är ettor.
  \item Nollställ stoppuret vid ett stopp och kör till nästa. Med \code{QUARTER_MS} = 1 ska två
    stopp i rad ligga ungefär 16~ms isär, 4 \enquote{sekunder} gånger 4~ms.
\end{enumerate}

I kursens kontroll av programmet, med \code{QUARTER_MS} = 250, tar en hel cykel
\textbf{128\,000\,716} cykler, 8,0000448~s. Med \code{QUARTER_MS} = 1 tar den 512\,716 cykler,
32~ms: ungefär 250 gånger fortare.

\subsection{Till kortet}
\label{c11:b7}
\begin{enumerate}
  \item \textbf{Sätt tillbaka \code{QUARTER_MS} till 250.} Det är det vanligaste felet i det här
    steget. Med 1 går hela cykeln runt ungefär 30 gånger i sekunden, och båda lamporna ser ut att
    lysa svagt på en gång.
  \item Koppla lysdioderna enligt \secref{c11:a5}, med USB-kabeln urkopplad.
  \item Bygg, anslut kortet och för över programmet enligt \secref{studio:5} i
    bilaga~\ref{app:studio}.
  \item Ta tid på några cykler och jämför med tabellen.
\end{enumerate}

Om lamporna gör något annat än i simulatorn, följ felsökningsordningen i \secref{c11:a6}.

\subsection{Från vägarbetsljus till trafikljus}
\label{c11:b8}
Slutuppgiftens \hyperref[lab3:a]{del~A}, trafikljuset i bilaga~\ref{app:lab3}, använder precis
samma metod med fler tillstånd:
\begin{itemize}
  \item \textbf{Tabellen} får fyra rader i stället för två, och en tredje lampa, gul på PB1.
  \item \textbf{Flödesplanen} får fyra par av symboler, ett per tillstånd.
  \item \textbf{Bitmönstren} blir fyra, och ett av dem har två lampor tända samtidigt.
  \item \textbf{Subrutinerna} är desamma, oförändrade.
\end{itemize}

\needspace{16\baselineskip}
Ljuscykeln som del~A ska ge, ritad som ett tidsdiagram:

\bookfigure{l11_traffic_timing}{Tidsdiagram för trafikljuset i slutuppgiftens del~A: en cykel på
8~s.}{c11:fig:traffic_timing}

\hyperref[lab3:b]{Del~B}, övergångsstället, lägger till något nytt: ett \textbf{villkor}. Ljuset
står still i ett tillstånd tills knappen trycks, och det blir en beslutssymbol i flödesplanen och
en väntan med \code{sbic} i programmet (\secref{c10:b5}).
